% Generated draft: requires mathematical and source review.
\documentclass[11pt]{article}
\usepackage[margin=25mm]{geometry}
\usepackage{fontspec}
\setmainfont{Latin Modern Roman}
\usepackage{amsmath,amssymb,mathtools,unicode-math}
\setmathfont{Latin Modern Math}
\usepackage{enumitem,needspace,xurl,hyperref}
\setlength{\parindent}{0pt}
\setlength{\parskip}{5pt}
\setlength{\emergencystretch}{3em}
\newcommand{\N}{\mathbb{N}}
\newcommand{\Z}{\mathbb{Z}}
\newcommand{\Q}{\mathbb{Q}}
\newcommand{\R}{\mathbb{R}}
\newcommand{\C}{\mathbb{C}}
\newcommand{\F}{\mathbb{F}}
\newcommand{\A}{\mathbb{A}}
\newcommand{\PP}{\mathbb{P}}
\newcommand{\E}{\mathbb{E}}
\newcommand{\eps}{\varepsilon}
\newcommand{\dd}{\,\mathrm{d}}
\begin{document}
\section*{UnsolvedMath ID: 30002534}
\textbf{Source problem code:} catalogue entry
\subsection{Definitions and mathematical statement}
Put \(D=K_X+\Delta\). A Kawamata log terminal (klt) pair consists, in the usual effective-boundary convention, of a normal variety \(X\) and an effective boundary divisor \(\Delta\), with \(D\) \(\R\)-Cartier, whose discrepancies are all strictly greater than \(-1\). The source assumes \(\kappa(X,\Delta)\geq 0\), understood as the nonnegativity of the Kodaira dimension of \(D\). For a \(\Q\)-Cartier divisor, this is the usual Iitaka dimension defined using integral Cartier multiples. The source does not specify the boundary coefficient class or a convention for Kodaira dimension beyond that setting.

In the absolute projective setting, let \(N_1(X)_{\R}\) be the real vector space of one-cycles modulo numerical equivalence, and let \(\overline{\mathrm{NE}}(X)\) be the closure of the cone generated by classes of integral curves. A ray \(R\subseteq\overline{\mathrm{NE}}(X)\) is extremal if \(\alpha+\beta\in R\), with \(\alpha,\beta\in\overline{\mathrm{NE}}(X)\), implies \(\alpha,\beta\in R\). It is \(D\)-negative if \(D\cdot\gamma<0\) for every nonzero \(\gamma\in R\).

\textbf{Source conjecture, with unresolved ambient hypotheses.} For every klt pair \((X,\Delta)\) satisfying \(\kappa(X,\Delta)\geq0\) to which the intended cone theorem applies, do there exist \(m\in\Z_{\geq0}\) and integral curves \(C_1,\ldots,C_m\subseteq X\) such that
\[
\bigl\{R\subseteq\overline{\mathrm{NE}}(X):R\text{ is a }D\text{-negative extremal ray}\bigr\}
=\bigl\{\R_{\geq0}[C_i]:1\leq i\leq m\bigr\}?
\]
Here the set on the right is empty when \(m=0\). Under the curve-generation conclusion of the intended cone theorem, this asks whether there are only finitely many \(D\)-negative extremal rays.

The projective notation above describes the absolute interpretation; the source does not explicitly specify projectivity, a ground field, characteristic, dimension, or whether a relative cone theorem is intended. These are unresolved hypotheses, so this is a conditional reconstruction of the stated target. The supplied background also mentions a restricted smooth threefold problem without stating its full hypotheses; it does not provide a precise replacement for the general conjecture.
\subsection{Short English statement}
The source asks whether a klt pair of nonnegative log Kodaira dimension has only finitely many negative extremal rays, each represented by one of finitely many curves. Its ambient hypotheses require clarification. Supplied dated curation labels the problem partially solved and reports failure for general klt pairs while mentioning restricted threefold cases; those claims are not independently verified here.
\subsection{Sources}
\begin{itemize}
\item Local registry: \texttt{data/unsolvedmath/problems.json}, supplied complete record for permanent UnsolvedMath ID \(30002534\), source code \texttt{catalogue entry}, entitled ``Finiteness of Negative Extremal Rays for KLT Pairs.'' The record belongs to the Oberwolfach Reports Open Problems collection and gives proposed year 2014; no individual proposer is supplied.
\item Supplied source citation: \emph{Mini-Workshop: Negative Curves on Algebraic Surfaces} (2014), Oberwolfach Reports, pp.~555--593, DOI 10.4171/owr/2014/10. Supplied literature label: ``Oberwolfach Report 10, 2014.'' \url{https://doi.org/10.4171/owr/2014/10}.
\item Supplied literature reference: ``Kapovich, 2013.'' \url{https://doi.org/10.1007/s00222-013-0453-4}. No title, precise result, or explanation of its applicability is supplied.
\item Literature assessment date: \texttt{2026-08-21}. This is supplied dated curation, not independently verified. It labels the problem partially solved, reports positive special cases in dimension three and failure for general klt pairs, and describes a relevant smooth threefold case as a remaining restricted problem. The supplied background explicitly requests independent verification of time-sensitive status and cited claims.
\item The registry labels the statement \texttt{corrected\_verified} and reports a manual semantic, typographic, and LaTeX review. This is supplied editorial metadata, not independent verification of the conjecture or its literature status.
\end{itemize}
\end{document}
